
\subsection{静电场}

\begin{tabularx}{\columnwidth}{XX}
  \multicolumn{2}{l}{\textbf{基本定义与等式}}\\
  场强定义&$E=\frac{F}{q}$\\
  电势定义&$\varphi = \frac{E_p}{q}$\\
  电势差定义&$U_{AB}=\varphi_A-\varphi_B$\\
  电场力做功&$W_{AB}=qU_{AB}=E_{pA}-E_{pB}$\\
  电势能变化&$\Delta E_p=-W_{\text{电}}$\\
  \multicolumn{2}{l}{\textbf{匀强电场规律}}\\
  电势差与场强关系&$U=Ed$（$d$ 为沿电场方向距离）\\
  \multicolumn{2}{l}{\textbf{点电荷}}\\
  库仑定律&$F=\frac{kQq}{r^2}$\\
  点电荷电场强度&$E=\frac{kQ}{r^2}$\\
  电场叠加&$\vec E=\vec E_1+\vec E_2+\cdots$\\
  \multicolumn{2}{l}{\textbf{电容器}}\\
  电容定义&$C=\frac{Q}{U}$\\
  电容决定式&$C=\varepsilon_0 \varepsilon_r \frac{S}{d}$\\
  真空介电常数关系&$k = \frac{1}{4\pi \varepsilon_0}$\\
  电容器能量&$E = \frac{1}{2} C U^2 = \frac{1}{2}\frac{Q^2}{C} = \frac{Q U}{2}$\\
  \multicolumn{2}{l}{\textbf{电场中的运动}}\\
  加速电场动能&$qU=\frac{1}{2}mv^2,\quad v=\sqrt{\frac{2qU}{m}}$\\
  加速电场动量&$qU=\frac{p^2}{2m},\quad p=\sqrt{2mqU}$\\
  竖直加速度&$a_y=\frac{qU}{md}$\\
\end{tabularx}

其中 $C$ 为电容，$\varepsilon_0$ 为真空介电常数，$\varepsilon_r$ 为介质的相对介电常数，$S$ 为极板面积，$d$ 为极板间距离。

\textbf{电容器描述：}
电容器存储电能的元件，由两个导体之间隔着绝缘材料构成。平行板电容器是由两块平行放置的导体板组成，中间夹有绝缘介质。常用于研究电容器的基本性质。

\textbf{电容器充放电：}
电容器的充电过程是指在外加电压作用下，电容器两极板之间积累电荷的过程；放电过程是指电容器释放储存的电能，电荷流向外部电路的过程。只有在电容器两极板间存在电势差（即外加电源或电压）时，电容器才能充电。充电过程中，电流方向为正极板流入、负极板流出，直到两极板间的电势差等于外加电压时，充电过程结束，电流为零。

\textbf{充电过程中电流和电压的变化：}
电容器刚开始充电时，极板间电压为零，电流最大；随着充电进行，极板间电压逐渐升高，电流逐渐减小；当极板间电压等于外加电压时，电流减小到零，充电过程结束。

电容器存在击穿电压，当电压过大时电容器中间的绝缘材料就会被击穿，使得电容器被破坏。

\textbf{电容器注意事项：}
电容器可以进行串并联，并联电容器，等效电容为电容器电容之和，串联电容器，等效电容规律和并联电阻相同。

\textbf{静电场注意事项：}
\begin{itemize}
  \item 电场力做功只与电荷初末位置有关，与路径无关
  \item 电势能增加，一般情况下，机械能减少
  \item 电场力做功时需要进行投影：$W=qEd\cos\theta$
  \item 电场线与等势面垂直
  \item 在匀强电场中，两点所在直线的场强均匀，可以通过$E'=\frac{U}{d}$求出
  \item 固定水平位移时，合速度最小时，水平与竖直方向速度相等
\end{itemize}

\begin{formulanotebox}
电场强度、电势、电势能分别对应“力学效应”“位置状态”“能量状态”，解题时不要混用单位和物理含义。

应用 $U=Ed$ 时需满足匀强电场且 $d$ 为沿场强方向的投影距离。
\end{formulanotebox}

\begin{keypointbox}
\textbf{电势与做功关系}：电场力做功只与初末位置有关，与路径无关；常用“先判电势高低，再判电势能变化”。

\textbf{电容器动态分析}：断开电源看 $Q$ 不变，接电源看 $U$ 不变，再结合 $C$ 的变化判断其他量变化。
\end{keypointbox}

\begin{casetypebox}[title={\strut\textcolor{white}{\bfseries 常见题型：电场做功与电容器动态分析}}]
\textbf{电场做功与电势能：}
\begin{itemize}[leftmargin=1.5em]
  \item 先确定电荷正负，再用 $W_{AB}=qU_{AB}$ 判断电场力做功正负。
  \item 电场力做正功时电势能减少，电场力做负功时电势能增加。
  \item 匀强电场中只取沿电场方向的位移，$U=Ed$。
\end{itemize}

\textbf{电容器动态分析：}
\begin{itemize}[leftmargin=1.5em]
  \item 接电源时 $U$ 不变，先由 $C=\varepsilon_0\varepsilon_r\frac{S}{d}$ 判断 $C$，再由 $Q=CU$ 判断 $Q$。
  \item 断开电源时 $Q$ 不变，先判断 $C$，再由 $U=\frac{Q}{C}$ 和 $E=\frac{U}{d}$ 判断电压、电场强度。
  \item 插入电介质或增大极板面积会使 $C$ 增大，增大板间距会使 $C$ 减小。
\end{itemize}
\end{casetypebox}
